% TOP_PROBLEM: {"id": 2648, "title": "Kourovka Notebook Problem 21.139", "category_id": 17, "category": {"id": 17, "name": "group_theory", "display_name": "Group Theory", "description": "Problems about groups, group actions, representations, and related algebraic structures.", "slug": "group-theory", "order_index": 17, "created_at": "2026-07-31T15:26:25.671Z"}, "status": "open", "problem_number": "KOU-21.139", "set_id": 10}
% FIRST_POSTED: 2026-10-01
% ENTRY_KIND: statement_only
% UnsolvedMath ID 2648; source catalogue code KOU-21.139
% !TeX program = lualatex
% Definition and sources only; this file is not an LLM solution attempt.
\documentclass[11pt,a4paper]{article}
\usepackage[margin=25mm]{geometry}
\usepackage{fontspec}
\setmainfont{lmroman10-regular.otf}[BoldFont=lmroman10-bold.otf,
 ItalicFont=lmroman10-italic.otf,BoldItalicFont=lmroman10-bolditalic.otf]
\usepackage{amsmath,amssymb,mathtools}
\usepackage{unicode-math}
\setmathfont{latinmodern-math.otf}
\AtBeginDocument{\let\setminus\smallsetminus}
\usepackage{xurl,hyperref}
\hypersetup{colorlinks=true,urlcolor=blue}
\setcounter{secnumdepth}{0}
\setlength{\parindent}{0pt}
\setlength{\parskip}{5pt}
\setlength{\emergencystretch}{3em}
\begin{document}
\section*{Kourovka Notebook Problem 21.139}\label{2648}
{\small UnsolvedMath ID 2648 · KOU-21.139 · Group Theory · Kourovka Notebook - New Problems, Issue 21}
\subsection{Definitions and mathematical statement}
A finitely generated group $G$ is word-hyperbolic if its Cayley graph for a
finite generating set has some $\delta\geq0$ such that every geodesic triangle
has each side in the $\delta$-neighborhood of the other two sides.
A right-angled Artin group has generators the vertices of a finite graph
$\Lambda$, with commuting relations exactly on its edges. Its Salvetti cube
complex has one vertex, one loop per generator, and a coordinate torus for
each complete subgraph, glued along the corresponding coordinate subtori.

Use the standard equivalent description of a compact special cube complex:
a compact nonpositively curved cube complex admitting a cubical local isometry
to a Salvetti complex. A cube complex glues unit Euclidean cubes along faces;
nonpositive curvature means every vertex link is flag (every finite clique
spans a simplex). A cubical local isometry induces an injective map of links
onto a full subcomplex, containing every simplex spanned by its vertices.
The group $G$ is virtually compact special if a
finite-index subgroup is the fundamental group of such a complex.
For $H\leq G$ let $b_2(H)=\dim_{\mathbb Q}H_2(H;\mathbb Q)$ be its second
rational group-homology Betti number, computed as the homology of a
classifying space $K(H,1)$ (a connected CW-complex with fundamental group
$H$ and contractible universal cover).

Suppose $G$ is hyperbolic and virtually compact special, and that
\[
\exists B<\infty\quad\forall H\leq G\text{ of finite index},\quad b_2(H)\leq B.
\]
Must $G$ have a finite-index subgroup that is a free group or the fundamental
group of a closed surface of negative Euler characteristic? Surfaces with
boundary have free fundamental groups and are included in the first alternative.
The bound must hold for all finite-index subgroups, not just one tower of covers.
\subsection{Short English statement}
Does uniformly bounded second homology in finite covers force a hyperbolic virtually compact special group to be virtually free or a surface group?
\subsection{Sources}
\begin{itemize}
\item[S1] ULAM.AI and the UnsolvedMath Contributors, supplied problems.json, record 2648 (KOU-21.139), Kourovka Notebook - New Problems, Issue 21. Catalogue identity and source transcription; CC BY 4.0.
\url{https://huggingface.co/datasets/ulamai/UnsolvedMath}.
\item[S2] The Kourovka Notebook, 21st edition, new problems, Problem 21.139. Expanded formulation of the supplied catalogue statement.
\url{https://arxiv.org/pdf/1401.0300v47}.
\end{itemize}
\end{document}
